%% kapa: exam=1 round=27 year=2016 period=1 form=A subject=civil_law subject_ko=민법 seq=6 number=6
%% source: Q-Net 감정평가사 기출문제 (cst00309) artlSeq=5169402 file=extracted/1008/5169402/제27회 감정평가사 1차 문제/제27회 감정평가사 1교시문제(A형).pdf
%% conversion: kordoc 4.24.51 -> tools/gen1.py (auto); review=proofread-sonnet; flags=none
\begin{question}{6}
사회질서에 반하는 법률행위에 관한 설명으로 옳지 않은 것은? (다툼이 있으면 판례에 따름)
\begin{choices}
  \item 오로지 보험사고를 가장하여 보험금을 취득할 목적으로 체결한 생명보험계약은 무효이다.
  \item 비자금을 소극적으로 은닉하기 위하여 임치한 것은 사회질서에 반하는 법률행위로 볼 수 없다.
  \item 부첩관계를 청산하면서 희생의 배상 내지 장래 생활대책 마련의 의미로 금원을 지급하기로 한 약정은 공서양속에 반하지 않는다.
  \item 반사회적 법률행위에 의한 무효를 가지고 선의의 제3자에게는 대항할 수 없다.
  \item 형사사건에 관한 변호사 성공보수 약정은 재판의 결과를 금전적 대가와 결부시키는 것으로서 사회질서에 위배되는 것으로 평가할 수 있다.
\end{choices}
\end{question}
\answer{4}
